% ==========================================================================
% Replacement for Section 1 (Introduction) of jsa/main.tex, i.e. the block
% currently at lines 107--195 ("\section{Introduction}" ... end of
% "What this paper does not claim").
%
% What changed and why -> REVAMP_NOTES.md.
% Numeric anchors are unchanged; they still resolve to results/ via audit_numbers.py.
% The one \todo{} marks the C2 experiment that does not exist yet.
% ==========================================================================

\section{Introduction}
\label{sec:intro}

Deploying a large language model on an 8\,GB device is no longer a research
curiosity; it is the default target for on-device assistants, and the entire
post-training quantisation (PTQ) literature is organised around it.

\paragraph*{The cache is not the second-order term.}
That literature spends its bits on the weights and treats the KV cache as either
free or as someone else's problem. At deployed context lengths this is
arithmetic rather than a hazard: from the published model cards a KV cache costs
112\,KiB per token for Qwen3-0.6B and 1.7B and 144\,KiB for 4B and 8B, so at 32k
tokens the cache is 3.76--4.83\,GB against 0.22--4.09\,GB of 4-bit weights, and
the two are equal at 1{,}919 tokens for 0.6B, 6{,}144 for 1.7B, 12{,}319 for 4B
and 21{,}388 for 8B. On an 8\,GB device the cache overtakes the weights long
before the context lengths these models are deployed at. The question is
therefore joint --- given a fixed byte budget, how many bits for the weights and
how many for the cache --- and it is the question this paper answers.

\paragraph*{The literature answers the two halves separately.}
One line optimises weight precision at a fixed width: rotations to flatten
activations \cite{quarot,spinquant,duquant,sun2024}, Hessian-aware rounding
\cite{gptq,frantar2022}, and mixed-precision schemes that spend more bits on
sensitive layers \cite{hawqv2,mixllm,gamma,alignquant,wang2018}. System-level
work co-designs the kernel with the format
\cite{lin2024,zhao2023,wei2024,guo2026,mixquant}, and low-rank compensation
recovers what rounding loses \cite{song2026,park2026,liu2025}. A second line
compresses the cache \cite{kivi,kvquant,mixkvq,kvpareto,duanmu2024,su2025,he2024}
and takes the weight precision as given. Serving stacks that manage the cache as
a first-class resource \cite{kwon2023} do not choose weights against cache
either. The nearest prior art, \cite{kvpareto}, does map a joint frontier over
cache bit-width, cache granularity and chunked prefill --- but its weight side is
fixed at 4-bit AWQ, so the split is not among its decision variables, its
frontier is indexed by memory consumed rather than by a per-bit accuracy cost
measured end to end, and it reports no usable-memory fraction for its hardware,
so the frontier cannot be re-targeted to a device with a different $\rhoSym$.

\paragraph*{Thesis: three inputs are assumed, and all three are measurable.}
The allocation literature assumes the objective (speed), the exchange rate
between the two sides (practice spends it on the weights), and the budget (the
datasheet number). We measure each one, and each measurement reverses the
assumption.

\emph{The objective.} On our stack 4-bit weights buy little or no time, because
dequantisation is not fused into the matmul; quantisation here purchases
capacity, not throughput, so latency cannot be the objective of the allocation.

\emph{The exchange rate.} With published implementations on both sides --- the
official GPTQ for weights, KIVI's own quantiser for the cache --- one cache bit
costs $9.0\times$ less accuracy than one weight bit at 0.6B and $13.3\times$ at
1.7B ($11.39\%$ against $1.27\%$ perplexity, $3.20\%$ against $0.24\%$). The
ordering reproduces on WikiText-2, and a 2-bit passkey experiment locates where
the claim stops: at 4 bits the scheme does not change long-context retrieval, at
2 bits it decides it.

\emph{The budget.} The fraction of nominal memory actually available to the
model differs $3.15\times$ between a discrete-GPU laptop ($\rhoSym=0.133$) and a
unified-memory edge module ($\rhoSym=0.419$), both sold as 8\,GB. Because the
cache grows with context while the weights do not, that difference changes the
price of a weight bit by $3.7\times$: dropping Qwen3-4B from 8-bit to 4-bit
weights multiplies the attainable context by $1.57\times$ on the laptop but by
$5.76\times$ on the module. A single published configuration cannot be optimal
on both.

\paragraph*{What we deliver.}
From those three inputs an allocator follows, and we state its status plainly:
it is a multiple-choice knapsack over measured per-device curves, solved by exact
dynamic programming, which is textbook. The contribution is not the optimiser.
It is the inputs the optimiser runs on --- a measured $\rhoSym$, a measured
workspace curve, and a per-bit accuracy cost measured on both sides of the same
model --- the finding that the objective has to be capacity rather than latency,
and the resulting policy, which we \emph{execute} on a shipped runtime rather
than model. On both devices that policy compresses the cache and leaves the
weights at 8 bits; on the laptop it is $10.84\%$ better in perplexity and $19\%$
faster per decode step than the deployed alternative at 4k, for $644$\,MB of
memory that both configurations fit in anyway.

\paragraph*{Contributions}
\begin{enumerate}
\item \textbf{The per-bit cost of the two sides runs against practice, and it is
      the quantity the allocator needs.} Measured with published implementations
      on both sides, one cache bit costs $9.0\times$ less accuracy than one
      weight bit at 0.6B and $13.3\times$ at 1.7B; the ordering reproduces on
      WikiText-2 ($10.9\times$ and $7.7\times$), so it is a property of the
      models and quantisers rather than of one corpus. A 2-bit passkey
      experiment bounds the claim: at 4 bits every scheme we tried recovers the
      key, at 2 bits the scheme decides (5/5 against 0/5)
      (Section~\ref{sec:results-bitcost}).

\item \textbf{The objective must be capacity, not latency.} On our stack 4-bit
      weights buy little or no time --- on 0.6B they cost time outright, because
      dequantisation is not fused into the matmul --- so an allocator that
      minimises latency optimises a quantity this hardware does not trade. We
      give the measured substitution frontier for both models
      (Section~\ref{sec:results-nospeed}).

\item \textbf{Capacity is device-dependent through a measurable fraction.} The
      usable-memory fraction is $0.133$ on a discrete-GPU laptop and $0.419$ on
      a unified-memory module, a $1.6\times$ difference in budget and a
      $3.15\times$ difference in the fraction lost; it moves the price of a
      weight bit by $3.7\times$. We report $\rhoSym$ as a measurement procedure
      rather than a constant, and report the module's \emph{measured} ceiling
      beside the cost model's prediction for it (Section~\ref{sec:rho}).

\item \textbf{The policy is executed, not modelled.} We solve the allocation
      exactly and run the configuration it recommends on a shipped runtime that
      stores the cache packed. On an 8\,GB laptop the recommendation is the
      difference between a length existing and not: a 4B model at 32k allocates
      9{,}033\,MiB with an fp16 cache and never completes, against 6{,}873\,MiB
      with a packed 8-bit cache, 52.6 tokens/s and $+0.01\%$ perplexity. The same
      quantiser then costs $26\times$ more accuracy on a smaller model of the same
      family --- the paper's thesis, appearing in the execution of its own policy
      (Section~\ref{sec:results-exec}, Table~\ref{tab:exec}).

\item \textbf{The stack sets the budget before any bit does.} A silent
      attention-dispatch failure costs $32\times$ in attention peak memory and
      $5.8\times$ in 4k prefill time, and removing it extends the attainable
      context $4\times$ before any quantisation decision is made. We release a
      model-free reproduction, because the effect leaves no trace in the
      framework's output (Section~\ref{sec:results-stack}).
\end{enumerate}

\paragraph*{Scope}
The measurements are made on two models (0.6B, 1.7B) whose quantiser families we
could match on both sides, on two devices, with the power mode fixed; we did not
instrument clocks, so a throttling explanation for a small difference cannot be
excluded from our data alone. We do not propose a new quantiser, and we do not
claim the ordering is model-independent --- the margin narrows as the model
grows, which is what a reviewer should expect if the mechanism is what we say it
is. We also report a hypothesis we tested and could not support: that the
optimal allocation reverses between prefill and decode
(Appendix~\ref{app:prereg}). Each of these bounds a specific claim; none of them
bounds the allocation result, which rests on measurements we state per device.
